\section{Introduction}
\label{sec:introduction}

Skin malignancies constitute the most frequently diagnosed cancers worldwide, with cutaneous melanoma responsible for the majority of skin-cancer mortality despite accounting for a minority of total lesions~\cite{Tschandl2018dataset,Esteva2017dermatologist}. Early recognition of melanoma and non-melanoma skin cancers—such as basal cell carcinoma and intraepithelial carcinoma—directly determines five-year survival and surgical morbidity. Yet diagnostic accuracy outside specialized dermatology clinics remains constrained by limited specialist access, high patient volume in primary care, and subtle morphological overlap between malignant lesions and benign pigmented mimics such as seborrheic keratoses or atypical melanocytic nevi~\cite{Tschandl2019comparison,Brinker2019deep}. Clinical decision support systems built on deep learning promise consistent first-line triage, provided they achieve calibrated probabilities, clinically defensible sensitivity, transparent visual explanations, and equitable accuracy across patient demographics~\cite{Esteva2019guide,OakdenRayner2020hidden}.

Over Workshops~1 through~5 of this course project, Group~8 engineered a reproducible machine learning stack for a radically different domain: character-level autoregressive sequence modeling on the \texttt{flwrlabs/shakespeare} federated benchmark~\cite{Caldas2018leaf}. That task involved 4.22~million overlapping 80-character dialogue windows spanning 1,129 theatrical characters across 36~plays, with an extreme 791,902:1 class imbalance across 76 discrete tokens. Solving it required strict entity-grouped data splits to prevent 98.75\% window-overlap leakage, mixed-precision training loops portable across AMD ROCm and CPU hardware, post-hoc logit adjustment for rare classes~\cite{Menon2020long}, reinforcement learning with a Double+Dueling Deep Q-Network for controllable generation, 1D residual variational autoencoders with Free Bits and cyclical KL annealing for anomaly detection, and 1D Wasserstein GANs with gradient penalties~\cite{Gulrajani2017improved} for controlled synthetic data augmentation.

Moving from discrete Renaissance English text to continuous two-dimensional dermatoscopic imaging raises a fundamental scientific question: which components of a domain-specific machine learning pipeline actually transfer across modalities, and where does transfer fail? Prior work demonstrated that even within computer vision, ImageNet pretraining offers far narrower feature-reuse benefits for medical imaging than conventional wisdom assumes~\cite{Raghu2019transfusion,Matsoukas2022what}. When the source domain is 1D symbolic text rather than natural photographs, naively forcing every previous model into healthcare would yield misleading engineering theater rather than scientific evidence. Honest negative transfer analysis—documenting precisely what fails and why—is therefore an explicit scientific objective of this workshop.

Here we execute a systematic transfer audit and clinical adaptation study on the public HAM10000 / ISIC~2018 Task~3 dermoscopy repository~\cite{Tschandl2018dataset,Codella2019skin}. First, we classify every model and infrastructure module from Workshops~1--5 into positive architectural/methodological transfers versus rejected or negative-control components. Second, we build a modality-specific preprocessing pipeline incorporating lesion-grouped splitting (\texttt{lesion\_id}), dermoscope vignette cropping, DullRazor morphological hair removal~\cite{Lee1997dullrazor}, Shades of Gray color constancy~\cite{Finlayson2004shades}, clinically valid dihedral augmentations, and DICOM interoperability with HIPAA Safe Harbor metadata auditing. Third, we implement a hybrid CNN-Transformer (\texttt{MedHybridTransformer}) directly importing the Pre-LN Transformer blocks of Workshop~2, alongside an ImageNet-pretrained \texttt{ResNet-18} visual transfer reference~\cite{He2016deep}, a 2D convolutional VAE adapted from Workshop~4, and a class-conditional 2D WGAN-GP adapted from Workshop~5, evaluated through four controlled ablation studies. Fourth, we evaluate every checkpoint using clinical operating points (sensitivity, specificity, PPV, NPV, temperature-scaled calibration~\cite{Guo2017calibration}, and Decision Curve Analysis~\cite{Vickers2006decision}), subgroup fairness audits~\cite{OakdenRayner2020hidden,Daneshjou2022disparities}, quantitative saliency alignment against ground-truth lesion masks via Grad-CAM~\cite{Selvaraju2017grad} and Attention Rollout~\cite{Abnar2020quantifying}, production drift simulations, and a regulatory and ethical deployment assessment grounded in FDA guidance~\cite{Food2021artificial,Gerke2020need}.
